\chapter{Instrukcja uruchomienia systemu}

Poniżej przedstawiono instrukcję uruchomienia poszczególnych modułów systemu \verb|Wydatkomat| w środowisku Linux.

\section{Backend (REST API)}

Uruchomienie backendu realizowane jest przy użyciu konteneryzacji Docker.

\begin{enumerate}
	\item Przejść do katalogu głównego projektu.
	\item W przypadku pierwszego uruchomienia skopiować plik konfiguracyjny:
	\begin{verbatim}
		cp .env.example .env
	\end{verbatim}
	\item Dostosować konfigurację w pliku \verb|.env| (np. dane do bazy danych, konfiguracja środowiska).
	\item Uruchomić aplikację przy użyciu Makefile:
	\begin{verbatim}
		make up
	\end{verbatim}
\end{enumerate}

Po wykonaniu powyższych kroków aplikacja backendowa zostanie uruchomiona wraz z wymaganymi usługami w kontenerach Docker.

\section{Aplikacja webowa}

Aplikacja webowa została przygotowana do uruchomienia w trybie produkcyjnym w środowisku lokalnym.

\begin{enumerate}
	\item Przejść do katalogu projektu:
	\begin{verbatim}
		cd expense-splitter-web
	\end{verbatim}
	\item Zainstalować zależności:
	\begin{verbatim}
		npm install
	\end{verbatim}
	\item Utworzyć plik konfiguracyjny:
	\begin{verbatim}
		cp .env.example .env
	\end{verbatim}
	\item Ustawić adres API w pliku \verb|.env|:
	\begin{verbatim}
		VITE_API_URL=https://wydatkomat.tech/api
	\end{verbatim}
	\item Zbudować aplikację:
	\begin{verbatim}
		npm run build
	\end{verbatim}
	\item Uruchomić podgląd wersji produkcyjnej:
	\begin{verbatim}
		npm run preview
	\end{verbatim}
\end{enumerate}

Aplikacja dostępna jest następnie pod adresem lokalnym wskazanym przez serwer Vite.

\section{Aplikacja mobilna}

Aplikacja mobilna została przygotowana w technologii React Native z wykorzystaniem Expo.

\begin{enumerate}
	\item Sklonować repozytorium:
	\begin{verbatim}
		git clone https://github.com/Rafal-wq/expense-splitter-mobile.git
	\end{verbatim}
	\item Przejść do katalogu projektu:
	\begin{verbatim}
		cd expense-splitter-mobile
	\end{verbatim}
	\item Zainstalować zależności:
	\begin{verbatim}
		npm install
	\end{verbatim}
	\item Utworzyć plik konfiguracyjny:
	\begin{verbatim}
		cp .env.example .env
	\end{verbatim}
	\item Ustawić adres backendu:
	\begin{verbatim}
		EXPO_PUBLIC_API_URL=https://wydatkomat.tech/api
	\end{verbatim}
	\item Uruchomić aplikację:
	\begin{verbatim}
		npx expo start
	\end{verbatim}
\end{enumerate}

Po uruchomieniu należy zeskanować kod QR przy użyciu aplikacji Expo Go.

\section{Aplikacja desktopowa}

Aplikacja desktopowa została napisana w języku Python i nie wymaga procesu budowania.

\begin{enumerate}
	\item Przejść do katalogu projektu:
	\begin{verbatim}
		cd expense-splitter-desktop
	\end{verbatim}
	\item Utworzyć wirtualne środowisko:
	\begin{verbatim}
		python3 -m venv venv
	\end{verbatim}
	\item Aktywować środowisko:
	\begin{verbatim}
		source venv/bin/activate
	\end{verbatim}
	\item Zainstalować zależności:
	\begin{verbatim}
		pip install -r requirements.txt
	\end{verbatim}
	\item Uruchomić aplikację:
	\begin{verbatim}
		python main.py
	\end{verbatim}
\end{enumerate}

Przy pierwszym uruchomieniu aplikacja wyświetla ekran logowania, a następnie inicjalizuje główne okno systemu.